\chapter{神经元不止一种器件}
% 完整版：chapters/18_eetypes.tex

在第一章里，我们按神经元释放的物质给它们分了类。工程师会换一种分法：看每个神经元当作什么器件来用。结果发现，相同的外表下藏着一整箱电子元件。

\section*{一箱电子元件}

第二章那只漏水的桶就是一个有泄漏的累加器：它把一个短窗口内的信号加起来。窗口非常小时，它就变成了符合检测器。只对变化作出响应、随即沉默的神经元，是一个让变化通过的滤波器。自己会工作的神经元是节拍器，也就是振荡器；如果它的节奏由输入来微调，那就是压控振荡器。发放一串串咔嗒的神经元是簇放电发生器，负责写出“划”。跟着声波节拍咔嗒的神经元是鉴相器。

\pencilfig{18}{\pencilart{18}}{相同的外表下藏着一整箱电子元件：电阻、电容、线圈、振荡器。}

还有一些器件我们以前没有提过。\emph{谐振器}像秋千：它对某个特定节奏的推动响应最好，就像秋千只有推得合拍才会越荡越高。在神经元内部，扮演弹簧的是一种抗拒变化的慢电流。\emph{无泄漏累加器}能长时间保持一个数值，就像一个不漏电的电容——当你向一侧看时，保持眼睛位置的网络就是这样工作的。但为此反馈必须调得几乎完美，这是一种代价高昂的调校。\emph{锁存器}是这样一种神经元：一次短推动就能让它进入长时间工作，直到被关掉为止，就像一个自锁开关；在控制肌肉的神经元里，这种模式由血清素开启。

\section*{变色龙器件}

这次分类最重要的发现是：它们不是不同的元件，而是同一个元件的不同\emph{模式}。同一个神经元白天可能是累加器，夜里却是簇放电发生器，全看电台告诉它什么。工程师会把大脑称为一块可重编程的逻辑阵列：电路只有一套，元件的用途由配置来改变。

神经元能够工作在所有这些模式下，这是测量到的。大脑如何利用模式切换来计算——大部分还是假说。

\fullversion{18}
